\setcounter{subsection}{5}% Preserve the source-visible paragraph number X.1.6 in the standalone class.
\subsection{}
\ifthenelse{\boolean{orig}}
{\label{X.I.6}}
{\label{X.1.6}}
\setcounter{proposition}{6}% The source paragraph occupies item 1.6 in the shared theorem sequence.
Supongamos todavía que~$f$ es propio y separable, pero no hagamos ya ninguna
hipótesis sobre~$f_*(\cal{O}_X)$, que corresponderá a un recubrimiento étale
~$Y'$ bien determinado de~$Y$, provisto además, sobre~$b$, del punto~$b'$
imagen de~$a$. Aplicando entonces~\Ref{X.1.4} al morfismo canónico~$X\to Y'$,
y suponiendo~$f$ sobreyectivo, la sucesión exacta~\Ref{X.1.4} se sustituye por
la siguiente, análoga a la sucesión exacta de homotopía de los espacios
fibrados en topología algebraica:
$$
\pi_1(\overline X_y,\overline a)\to \pi_1(X,a)\to \pi_1(Y,b)\to
\pi_0(\overline X_y,\overline a)\to \pi_0(X,a)\to \pi_0(Y,b)\to e.
$$
Por supuesto, en~\Ref{X.1.4} no puede afirmarse en general que el
homomorfismo~$\pi_1(\overline X_y,\overline a)\to \pi_1(X,a)$ sea inyectivo;
en topología algebraica su núcleo es la imagen de~$\pi_2(Y,b)$, y también en
geometría algebraica habría que introducir grupos de homotopía en todas las
dimensiones y la sucesión exacta de homotopía completa para un morfismo propio
que satisfaga hipótesis convenientes (por ejemplo, ser un morfismo
\ifthenelse{\boolean{orig}}
{liso.}
{liso).}
En la actualidad no se dispone de ningún resultado en este sentido, salvo una
definición razonable (si no definitiva) de los grupos de homotopía superiores.

\begin{corollaire}
\label{X.1.7}
Sean~$k$ un cuerpo algebraicamente cerrado y~$X$ e~$Y$ dos preesquemas conexos
sobre~$k$; supongamos~$X$ propio sobre~$k$ e~$Y$ localmente
\marginpar{265}
noetheriano. Sean~$a$ un punto geométrico de~$X$ y~$b$ un punto geométrico
de~$Y$, con valores en una misma extensión algebraicamente cerrada~$K$ de~$k$;
consideremos el punto geométrico~$c=(a,b)$ de~$X\times_k Y$ y el homomorfismo
~$\pi_1(X\times_k Y,c)\to \pi_1(X,a)\times\pi_1(Y,b)$ deducido de los
homomorfismos de grupos fundamentales asociados a las dos proyecciones
~$X\times_k Y\to X$ y~$X\times_k Y\to Y$. El homomorfismo precedente es un
\emph{isomorfismo}.
\end{corollaire}

Supongamos primero~$K=k$. Pongamos~$Z=X\times_k Y$, consideremos la proyección
~$f\colon Z\to Y$ y el punto subyacente~$y$ del punto geométrico~$b$ de~$Y$,
y apliquemos a esta situación el resultado~\Ref{X.1.4}. Obsérvese para ello
que, tras pasar a~$X_\red$ (lo que no cambia los grupos fundamentales
considerados), puede suponerse~$X$ reducido y, por tanto, separable sobre~$k$;
así,
\ifthenelse{\boolean{orig}}
{$Z$ es separable sobre~$k$}
{$Z$ es separable sobre~$Y$},
y tiene evidentemente fibras geométricamente conexas (pues~$X$ es conexo). La
fibra geométrica de~$Z$ en~$b$ es canónicamente isomorfa a
~$X\otimes_k K=X$. Por otra parte, como el compuesto de los morfismos
~$X\to Z\to X$ es la identidad, se ve que
~$\pi_1(X,a)\to \pi_1(Z,c)$ es inyectivo, y~\Ref{X.1.4} nos da una sucesión
exacta:
$$
e\to \pi_1(X,a)\to \pi_1(Z,c)\to \pi_1(Y,b)\to e.
$$
Por otra parte, se tiene la sucesión exacta canónica:
$$
e\to \pi_1(X,a)\to \pi_1(X,a)\times \pi_1(Y,b)\to \pi_1(Y,b)\to e,
$$
donde los dos homomorfismos escritos son la inyección canónica y la proyección
canónica. Por último, se tiene un homomorfismo de la primera sucesión exacta en
la segunda mediante los morfismos identidad en los términos extremos y el
homomorfismo canónico~$\pi_1(Z,c)\to \pi_1(X,a)\times\pi_1(Y,b)$ en los
términos medios. La conmutatividad del diagrama así obtenido se comprueba de
inmediato. Como los homomorfismos en los términos extremos son isomorfismos,
ocurre lo mismo con los términos medios, lo que prueba~\Ref{X.1.7} en este
caso.

Cuando ya no se supone~$K=k$, solo se obtiene un isomorfismo
$$
\pi_1(Z,c)\to \pi_1(X\otimes_k K,a)\times \pi_1(Y,b),
$$
\marginpar{266}%
y~\Ref{X.1.7} equivale entonces al caso particular siguiente:

\begin{corollaire}
\label{X.1.8}
Sean~$X$ un esquema propio y conexo sobre un cuerpo algebraicamente cerrado
~$k$,~$k'$ una extensión algebraicamente cerrada de~$k$,~$a'$ un punto
geométrico de~$X\otimes_k k'$ y~$a$ su imagen en~$X$. Entonces el
homomorfismo canónico~$\pi_1(X\otimes_k k',a')\to \pi_1(X,a)$ es un
\emph{isomorfismo}.
\end{corollaire}

Que este homomorfismo sea sobreyectivo equivale a decir que, si~$X'$ es un
recubrimiento étale conexo de~$X$, entonces~$X'\otimes_k k'$ es también
conexo, y resulta inmediatamente del hecho de que~$k$ es algebraicamente
cerrado; es también un caso particular
\ifthenelse{\boolean{orig}}
{\Ref{IX.3.4}.}
{de~\Ref{IX.3.4}.}
La hipótesis de que~$X$ sea propio no ha intervenido todavía. Dicho esto, la
inyectividad del homomorfismo considerado significa también lo siguiente:
\emph{todo recubrimiento étale de~$X\otimes_k k'$ es isomorfo a la imagen
inversa de un recubrimiento étale de~$X$.} Es esencialmente un sorites comprobar
que puede encontrarse una subálgebra sobre~$k$,~$A$, de
\ifthenelse{\boolean{orig}}
{$K$}
{$k'$},
de tipo finito sobre~$k$, y un recubrimiento étale de~$X\otimes_k A$ cuya
imagen inversa sobre~$X\otimes_k k'$ sea isomorfa al recubrimiento dado. Sea,
pues,~$Y=\Spec(A)$, que es un~$k$-esquema íntegro de tipo finito y tiene, por
tanto, puntos racionales sobre~$k$. Apliquemos entonces~\Ref{X.1.7} al grupo
fundamental de~$X\times Y$ en un punto~$(a,b)$ racional sobre~$k$: se obtiene
que todo recubrimiento étale conexo de~$X\times Y$ es isomorfo a un cociente
de un recubrimiento~$X'\times Y'$, donde~$X'$ e~$Y'$ son recubrimientos
galoisianos étale de~$X$ y~$Y$, de grupos~$G$ y~$G'$, por un subgrupo~$H$
de~$G\times G'$. Esto implica que la imagen inversa de este recubrimiento de
~$X\times Y$ sobre~$X\times Y'$ es isomorfa a un recubrimiento de la forma
~$X_1'\times Y'$, donde~$X_1'$ es un recubrimiento étale de~$X$. Si~$L$ es
el cuerpo de funciones de~$Y$, igual al cuerpo de fracciones de~$A$ en~$k'$,
el recubrimiento étale de~$X\otimes_k L$ inducido por el recubrimiento dado de
~$X\times_k Y$ es tal que existe una extensión finita separable~$L'$ de~$L$
para la cual la imagen inversa de dicho recubrimiento sobre~$X\otimes_k L'$
es isomorfa a~$X_1'\otimes_k L'$. Ahora bien, como~$k'$ es algebraicamente
cerrado, puede suponerse que la extensión~$L'$ de~$L$ está contenida en~$k'$.
Esto prueba que el recubrimiento étale dado de~$X\otimes_k k'$ es isomorfo a
~$X_1'\otimes_k k'$, c.q.d.

La forma explícita señalada de paso para los recubrimientos étale de un
producto
\marginpar{267}
~$X\times_k Y$ implica de inmediato el resultado siguiente:

\begin{corollaire}
\label{X.1.9}
Sean~$k$ un cuerpo algebraicamente cerrado,~$X$ e~$Y$ dos preesquemas
localmente noetherianos sobre~$k$,~$Z=X\times_k Y$ su producto y~$Z'$ un
recubrimiento étale de~$Z$. Para todo punto~$y\in Y$ racional sobre~$k$, sea
~$i_y \colon \Spec(k)\to Y$ el morfismo canónicamente asociado,
~$j_y=\id_X\times_k i_y$ el morfismo correspondiente~$X\to Z$. Sea, por
último,~$X_y'$ el recubrimiento étale de~$X$ imagen inversa de~$Z'$ por~$j_y$.
Supongamos~$Y$ conexo y~$X$ o~$Y$ propio sobre~$k$. Entonces los
recubrimientos~$X_y'$ de~$X$ son todos isomorfos.
\end{corollaire}

De manera figurada, puede decirse que \emph{una familia de recubrimientos
étale de~$X$, parametrizada por un preesquema conexo~$Y$, es constante si~$X$
o el preesquema de parámetros~$Y$ es propio sobre~$k$.}

\begin{remarques}
\label{X.1.10}
Los corolarios~\Ref{X.1.7} a~\Ref{X.1.9} se deben a
Lang--Serre~\cite{X.2} en el caso de los esquemas algebraicos normales. (Su
trabajo fue la motivación inicial de la teoría del grupo fundamental
desarrollada en este Seminario). Como observaron estos autores, estos
resultados dejan de ser ciertos cuando se abandona la hipótesis de que sean
propios, al menos en característica~$p>0$. Tomando, por ejemplo, como~$X$ la
recta afín~$X=\Spec(k[t])$, no es difícil ver que los recubrimientos de~$X$,
parametrizados por la recta afín~$Y=\Spec(k[s])$, definidos por las ecuaciones
$$
x^p - x = st,
$$
son étale y no isomorfos dos a dos. Esto contradice~\Ref{X.1.9} y, con mayor
razón,~\Ref{X.1.7}; se ve igualmente que, si~$s$ se considera como un elemento
trascendente sobre~$k$ en una extensión algebraicamente cerrada~$K$ de~$k$,
se obtiene un recubrimiento étale~$X'$ de~$X\otimes_kK$ que no procede de un
recubrimiento étale de~$X$.
\end{remarques}
